% Relatório autossuficiente: não requer figuras ou arquivos de dados externos.
% Compilar duas vezes com pdfLaTeX para atualizar as referências internas.
\documentclass[11pt,a4paper]{article}
\usepackage[utf8]{inputenc}
\usepackage[T1]{fontenc}
\usepackage[brazil]{babel}
\usepackage[margin=2.3cm]{geometry}
\usepackage{booktabs,tabularx,array}
\usepackage{xurl}
\usepackage[hidelinks]{hyperref}
\urlstyle{tt}
\setlength{\parindent}{0pt}
\setlength{\parskip}{0.55em}
\setlength{\emergencystretch}{2em}
\renewcommand{\arraystretch}{1.15}
\newcommand{\hex}[1]{\texttt{0x#1}}

\title{Relatório de simulação RTL e cobertura\\
RSA de 32 bits com interface AMBA APB\\
\large Xcelium/Xrun --- clock de 100 MHz}
\author{Projeto RSA\_APB\_32-bit}
\date{Execução de 5 de novembro de 2026, às 15h10min39s (UTC$-$3)}
\begin{document}
\maketitle

\begin{abstract}
A última execução RTL disponível do periférico RSA/APB concluiu as 398
verificações previstas, com 398 acertos e nenhuma falha. O teste utilizou
clock de 100 MHz, 30 rodadas de banco de registradores e dez operações RSA,
com duração simulada de 134,98~$\mu$s. O IMC registrou, na hierarquia do DUT,
80,62\% de cobertura agregada, 80,45\% de cobertura de código e 90,12\%
de cobertura de FSM, nas colunas \emph{Covered}. Não foram registradas
mensagens de erro ou fatal. Há avisos sobre o escopo da instrumentação de
cobertura, e a cobertura funcional aparece como não disponível.
\end{abstract}

\section{Identificação e condições da execução}
Este relatório analisa a execução abaixo, sem executar nova simulação.
Ela substitui, como execução RTL mais recente, o teste das 15h07min12s.
O log confirma que o arquivo compilado nesta execução foi
\path{testbenchs/rsa_apb_top/testbench.sv}; o arquivo
\path{testbench_2.sv} não integra a linha de comando desta execução.

\textbf{Identificador da execução:}\\
\path{amba_3_apb_rsa_apb_top_2026-11-05_15-10-39}

\begin{center}
\begin{tabularx}{\linewidth}{@{}lX@{}}
\toprule
Item & Configuração registrada \\
\midrule
Topo RTL & \path{amba_3_apb_rsa_apb_top} \\
Instância do DUT & \path{testbench.dut} \\
Simulador & Cadence Xcelium/Xrun 24.03-s013 \\
Análise de cobertura & Cadence IMC 24.03-s002 \\
Script & \path{scripts/amba_3_apb_rsa_apb_top/bash/xrun.sh} \\
Fontes & RTL APB em \path{src/amba_3_apb/} e núcleo em \path{src/rsa32/v2/} \\
Parâmetros do testbench & Dados: 32 bits; endereço: 12 bits; \texttt{LOAD\_GAP=1} \\
Clock & 10,0 ns / 100 MHz \\
Escala temporal & 1 ns / 1 ps \\
Instrumentação & \texttt{-coverage all}; acesso \texttt{+rwc} \\
Modelo de simulação & RTL, sem netlist de células ou anotação SDF \\
\bottomrule
\end{tabularx}
\end{center}

Esta execução é independente de PDK. O clock é produzido pelo testbench;
o SDC não é uma entrada desta simulação. O funcionamento a 100 MHz no
modelo RTL não constitui validação dos atrasos físicos para essa frequência.

\section{Método de verificação}
O testbench aplica transações APB e compara os valores observados com um
scoreboard. Os sinais de entrada APB são alterados na borda de descida do
clock. A instância RTL recebe explicitamente \texttt{LOAD\_GAP=1}, que
introduz intervalo entre cargas dos operandos do núcleo RSA.

O plano executado contém oito verificações iniciais, 30 rodadas de banco
de registradores com oito verificações por rodada e dez rodadas RSA com
15 verificações por rodada:
\[
N_{\mathrm{checks}}=8+30\times8+10\times15=398.
\]
As operações RSA incluem carga e preservação de operandos, estados
\texttt{BUSY}/\texttt{DONE}, resultado, interrupção e retorno ao estado
ocioso. A referência aritmética é a função de exponenciação modular do
testbench. O limite de uma transação APB é de 20 ciclos; cada operação
RSA admite até 10\,000 consultas de estado. Há também um timeout global.

\section{Resultados funcionais}
\begin{center}
\begin{tabular}{@{}lr@{}}
\toprule
Métrica & Resultado \\
\midrule
Verificações previstas / realizadas & 398 / 398 \\
Acertos / falhas & 398 / 0 \\
Rodadas do banco de registradores / RSA & 30 / 10 \\
Tempo simulado & 134\,980 ns (134,98 $\mu$s) \\
Tempo de execução informado pelo Xrun & 2 s \\
Mensagens \texttt{*E} / \texttt{*F} no log Xrun & 0 / 0 \\
FSMs extraídas para cobertura & 4 \\
\bottomrule
\end{tabular}
\end{center}

O log registra \texttt{TEST PASS: 398 checks, 0 errors}, seguido de
encerramento normal por \texttt{\$finish}. A Tabela~\ref{tab:vectors}
reproduz os resultados observados nas dez rodadas RSA. Mensagens, módulos
e resultados estão em hexadecimal; rodada e expoente estão em decimal.

\begin{table}[htbp]
\centering
\small
\begin{tabular}{@{}rllrl l@{}}
\toprule
Rodada & Mensagem & Módulo & Expoente & Esperado = lido & Estado \\
\midrule
1  & \hex{00000041} & \hex{00000ca1} & 17 & \hex{00000ae6} & PASS \\
2  & \hex{0000007b} & \hex{00000ca1} & 7  & \hex{00000b34} & PASS \\
3  & \hex{00000000} & \hex{00000061} & 5  & \hex{00000000} & PASS \\
4  & \hex{00000001} & \hex{00000101} & 13 & \hex{00000001} & PASS \\
5  & \hex{0000002a} & \hex{00000061} & 0  & \hex{00000001} & PASS \\
6  & \hex{12345678} & \hex{0000fff1} & 1  & \hex{00006793} & PASS \\
7  & \hex{80000001} & \hex{fffffffb} & 2  & \hex{4000000b} & PASS \\
8  & \hex{ffffffff} & \hex{fffffffb} & 3  & \hex{00000040} & PASS \\
9  & \hex{075bcd15} & \hex{00010001} & 7  & \hex{00008a8e} & PASS \\
10 & \hex{fedcba98} & \hex{8000000b} & 17 & \hex{14adadfc} & PASS \\
\bottomrule
\end{tabular}
\caption{Vetores RSA e resultados observados no log da simulação.}
\label{tab:vectors}
\end{table}

O resultado da décima operação, \hex{14adadfc}, foi verificado em
134\,830 ns. As últimas verificações confirmaram a preservação dos
operandos, \texttt{STATUS.IDLE}=0 e IRQ=0 em 134\,980 ns. O percentual de
acertos do scoreboard é 100\%; ele é distinto das métricas de cobertura.

\section{Resultados de cobertura}
O IMC carregou o banco desta execução e produziu o relatório por instância
com \texttt{report -summary -cumulative on -inst -local off}.
A Tabela~\ref{tab:coverage} mantém a distinção entre \emph{Average}
e \emph{Covered} usada pelo relatório. Os percentuais \emph{Covered}
correspondem às razões de itens cobertos sobre itens contabilizados.
\emph{Average} é a média apresentada pelo IMC e não deve ser interpretada
como essa mesma razão.

\begin{table}[htbp]
\centering
\begin{tabular}{@{}llrr@{}}
\toprule
Escopo & Métrica & Average & Covered \\
\midrule
Testbench e descendentes & Overall & 69,75\% & 80,39\% (3836/4772) \\
Testbench e descendentes & Code & 69,69\% & 80,22\% (3763/4691) \\
Testbench e descendentes & FSM & 90,49\% & 90,12\% (73/81) \\
\midrule
DUT e descendentes & Overall & 74,36\% & 80,62\% (3616/4485) \\
DUT e descendentes & Code & 74,24\% & 80,45\% (3543/4404) \\
DUT e descendentes & FSM & 90,49\% & 90,12\% (73/81) \\
\bottomrule
\end{tabular}
\caption{Cobertura cumulativa por hierarquia, conforme relatório IMC.}
\label{tab:coverage}
\end{table}

A cobertura \emph{Functional} aparece como \texttt{n/a} em todas as
instâncias. A opção \texttt{-coverage all} não demonstra, por si só,
a existência de um modelo de cobertura funcional. Neste relatório,
``cumulativa'' inclui descendentes da instância; não indica união de
execuções distintas. O IMC carregou apenas o banco identificado na Seção 1.

\begin{center}
\small
\begin{tabular}{@{}lrr@{}}
\toprule
Instância (nome relativo) & Overall Covered & FSM Covered \\
\midrule
\texttt{data\_io\_top\_inst} & 59,40\% (354/596) & n/a \\
\quad\texttt{adrec\_inst} & 35,29\% (6/17) & n/a \\
\quad\texttt{da\_o\_m} & 59,26\% (160/270) & n/a \\
\texttt{apb\_control\_inst} & 76,12\% (51/67) & n/a \\
\texttt{regbank\_inst} & 69,36\% (335/483) & n/a \\
\texttt{reset\_inst} & 60,00\% (6/10) & n/a \\
\texttt{rsa\_wrapper\_inst} & 89,31\% (2649/2966) & 90,12\% (73/81) \\
\quad\texttt{rsa\_wrapper\_control\_inst} & 83,13\% (69/83) & 87,50\% (14/16) \\
\quad\texttt{rsa\_wrapper\_data\_inst} & 99,33\% (148/149) & n/a \\
\quad\texttt{rsa\_core\_inst} & 89,54\% (2259/2523) & 90,77\% (59/65) \\
\quad\texttt{rsa\_core\_mult\_blk} & 90,87\% (637/701) & 100,00\% (13/13) \\
\quad\texttt{rsa\_core\_mod\_blk} & 98,36\% (899/914) & 85,00\% (17/20) \\
\quad\texttt{rsa\_core\_ctrl\_blk} & 74,65\% (480/643) & 90,63\% (29/32) \\
\bottomrule
\end{tabular}
\end{center}

As linhas são cumulativas e se sobrepõem hierarquicamente; seus numeradores
não devem ser somados. Restam 869 itens não cobertos no agregado do DUT,
sendo 861 de código e oito de FSM. Os oito itens FSM distribuem-se entre
controle do wrapper (dois), bloco modular (três) e controle do núcleo (três).
O resumo não identifica quais estados ou transições faltam.

O decodificador de endereços apresenta a menor cobertura agregada entre
os blocos selecionados. Isso orienta uma inspeção detalhada dos itens
não cobertos antes de definir novos estímulos ou exclusões. Não se pode
concluir, apenas pelo resumo, se cada item representa comportamento
alcançável, lógica constante ou cenário não exercitado.

\section{Avisos e limites da instrumentação}
O Xrun emitiu seis avisos, um de cada código abaixo. O IMC emitiu um aviso
\texttt{REPDEP}. Não foram encontradas mensagens de erro ou fatal nos
logs dessas duas ferramentas analisados.

\begin{center}
\begin{tabularx}{\linewidth}{@{}lX@{}}
\toprule
Código & Interpretação e impacto \\
\midrule
\texttt{COVSCE} & Cobertura de expressões com cast estático habilitada por padrão; atenção à compatibilidade ao unir bancos antigos. \\
\texttt{COVFDP} & Comparação de descrições FSM entre instâncias transferida para o descarregamento dos dados de simulação; aviso sobre configuração de módulos parametrizados. \\
\texttt{COVFHT} & Transições de permanência no mesmo estado não são extraídas por padrão. \\
\texttt{COVDCL} & A cobertura de expressões padrão limita operadores e contextos instrumentados; não inclui automaticamente todas as expressões RTL. \\
\texttt{COVMDD} & Toggle de arrays estáticos multidimensionais não habilitado por padrão, com referência ao somador Kogge--Stone. \\
\texttt{COVUTA} & Tipos como strings e certas estruturas de arrays não suportam toggle; referência a \texttt{failure\_details} no testbench. \\
\texttt{REPDEP} & Comando de relatório legado do IMC; a ferramenta recomenda \texttt{report\_metrics} para relatórios dinâmicos. \\
\bottomrule
\end{tabularx}
\end{center}

Esses avisos não indicam falha das comparações do scoreboard. Eles delimitam
o universo de cobertura medido: mesmo 100\% de uma categoria não abrangeria
itens excluídos da instrumentação. O relatório resumido também não separa
as contagens de bloco, expressão e toggle dentro de \emph{Code}.

\section{Comparação com a simulação com SDF}
A execução pós-síntese de referência é
\path{amba_3_apb_rsa_apb_top_2026-11-05_14-37-30_767454205_tsmc28},
documentada em \path{docs/relatorio_simulacao_sdf_tsmc28_2026-11-05.tex}.
A comparação textual das 398 linhas numeradas dos scoreboards encontrou
igualdade de valores esperados, valores lidos, status e tempos de verificação.
Ambas as execuções concluíram em 134\,980 ns.

Essa concordância confirma consistência nos pontos observados pelos testes.
Ela não equivale a prova formal de equivalência nem demonstra identidade
de todas as transições internas. Na execução RTL não há anotação de atrasos
de células; as limitações de timing descritas no relatório SDF permanecem
específicas daquela execução. Não há métrica IMC da execução SDF apresentada
ali que permita uma comparação direta de cobertura com estes valores.

\section{Conclusão}
A regressão RTL a 100 MHz passou nas 398 verificações, com os dez resultados
RSA corretos e concordância com o scoreboard da execução pós-síntese.
A cobertura medida no DUT alcançou 80,62\% no agregado e 90,12\% em FSM.
O teste fornece evidência funcional para os cenários executados, mas ainda
há itens de código e FSM não cobertos e não há cobertura funcional reportada.

Os próximos passos sugeridos são inspecionar os itens não cobertos por
instância, revisar o escopo de instrumentação indicado pelos avisos e
complementar os testes conforme o plano funcional. Nenhuma alteração de
RTL, testbench ou scripts foi necessária para elaborar este relatório.

\appendix
\section{Rastreabilidade dos artefatos}
Nos caminhos abaixo, \texttt{RUN} representa o identificador completo da
Seção 1. Todos os caminhos são relativos à raiz do repositório.
\begin{itemize}
\item Log Xrun: \path{logs/amba_3_apb_rsa_apb_top/xrun/<RUN>.log}.
\item Relatório IMC: \path{reports/amba_3_apb_rsa_apb_top/<RUN>_coverage.rpt}.
\item Log do IMC: \path{logs/amba_3_apb_rsa_apb_top/imc/imc.log}.
Este arquivo não possui identificador no nome; o conteúdo analisado registra
explicitamente o carregamento da execução das 15h10min39s.
\item Banco de cobertura: \path{coverage/amba_3_apb_rsa_apb_top/scope/<RUN>/icc_0cdf8577_00000000.ucd}.
\item Modelo de cobertura: \path{coverage/amba_3_apb_rsa_apb_top/scope/icc_0cdf8577_00000000.ucm}.
\item Formas de onda: \path{waveforms/amba_3_apb_rsa_apb_top/<RUN>.vcd}.
O arquivo existe e possui 79\,412\,636 bytes; seu conteúdo não foi necessário
para as conclusões baseadas nos logs e relatórios.
\item Testbench: \path{testbenchs/rsa_apb_top/testbench.sv}.
\item Script: \path{scripts/amba_3_apb_rsa_apb_top/bash/xrun.sh}.
\end{itemize}
\end{document}
